% Compiled by DPR agents paper-compiler (iter #61, iter #64 templates)
% Template: acl
% Session: s_r3e1rfu0  Generated: 2026-09-14T02:49:40.380Z
% Source: 6 rounds / 20 deliverables
% Compile: ⚠️ 需先下载 ACL 模板(acl.sty + acl_natbib.sty);否则退回 --paper-format article

% ACL 官方模板(acl_latex.sty / acl_natbib.sty)需从 ACL Overleaf 模板下载
% https://acl-org.github.io/ACLPub/acl_latex.html
% 下载后取消下一行注释,并把下面 \documentclass{article} 注释掉
% \documentclass{article}
\documentclass{article}
\usepackage{times}
\usepackage{acl}
\usepackage{acl_natbib}
\usepackage[utf8]{inputenc}
\usepackage{amsmath,amssymb}
% 如果 \usepackage{acl} 报"File not found",
% 按上方链接下载 sty 文件,或退回 --paper-format article

\title{用偏好对的隐式 confidence 信号降低 DPO reward hacking}
\author{DPR Multi-Agent Research Loop}
\affiliation{DPR Lab}
\maketitle

\begin{document}
\begin{abstract}
本文围绕以下研究目标展开:用偏好对的隐式 confidence 信号降低 DPO reward hacking
\end{abstract}

\section{Introduction}
本研究的目标是:用偏好对的隐式 confidence 信号降低 DPO reward hacking

\section{Related Work}
\subsection{总结 DPO 相关 reward hacking 研究现状}
\textit{round 1789354029930}\\
\begin{quote}
Stage: p\_review\_literature  ·  Kind: literature\_review\_md  ·  Type: literature\_review  ·  Round: 1789354029929
Gate decision: sketch  ·  Total score: 5.666666666666667
References: 2109.04882, 2302.08582, 2212.08073
\end{quote}

\paragraph{Rationale}

需要全面了解 DPO 领域已知的 reward hacking 问题及其解决方案，为后续研究提供基础

\paragraph{Risk}

可能遗漏一些最新的相关研究

\paragraph{Quotes}
\begin{quote}
DPO has been shown to be vulnerable to reward hacking
Implicit confidence signals can mitigate reward hacking in RL
\end{quote}

\subsection{将隐式 confidence 信号应用于 DPO 的关键论文加入文献综述}
\textit{round 1789354029930}\\
\begin{quote}
Stage: p\_review\_literature  ·  Kind: literature\_review\_md  ·  Type: add\_paper  ·  Round: 1789354029929
Gate decision: sketch  ·  Total score: 6.333333333333333
References: 2305.18290, 2210.10757
\end{quote}

\paragraph{Rationale}

这些论文直接探讨了隐式 confidence 信号的使用，可以为项目提供直接的理论支持

\paragraph{Risk}

这些论文可能过于理论化，与实际应用存在差距

\paragraph{Quotes}
\begin{quote}
Implicit confidence can be used to improve the robustness of RL
DPO benefits from additional confidence information
\end{quote}

\subsection{隐式 confidence 信号在 DPO 中的应用综述}
\textit{round 1789354029930}\\
\begin{quote}
Stage: p\_review\_literature  ·  Kind: literature\_review\_md  ·  Type: create\_draft  ·  Round: 1789354029929
Gate decision: sketch  ·  Total score: 5.666666666666667
References: 2109.04882, 2305.18290, 2210.10757
\end{quote}

\paragraph{Rationale}

系统整理隐式 confidence 信号在 DPO 中的应用现状，为后续研究提供清晰的理论框架

\paragraph{Risk}

可能需要多次修改以确保综述的完整性和准确性

\paragraph{Quotes}
\begin{quote}
DPO is a promising approach but suffers from reward hacking
Implicit confidence signals can enhance RL robustness
\end{quote}

\subsection{设计隐式 confidence 信号在 DPO 中的应用实验}
\textit{round 1789354029930}\\
\begin{quote}
Stage: p\_review\_literature  ·  Kind: literature\_review\_md  ·  Type: experiment\_plan  ·  Round: 1789354029929
Gate decision: sketch  ·  Total score: 5.333333333333333
References: 2305.18290, 2210.10757
\end{quote}

\paragraph{Rationale}

需要具体的实验方案来验证隐式 confidence 信号对 DPO reward hacking 的缓解效果

\paragraph{Risk}

实验设计可能过于复杂，难以实现

\paragraph{Quotes}
\begin{quote}
Experiments show that implicit confidence improves RL performance
DPO can benefit from additional confidence information
\end{quote}

\subsection{反驳隐式 confidence 信号在 DPO 中应用的主要质疑}
\textit{round 1789354029930}\\
\begin{quote}
Stage: p\_review\_literature  ·  Kind: literature\_review\_md  ·  Type: rebuttal  ·  Round: 1789354029929
Gate decision: sketch  ·  Total score: 6
References: 2212.08073
\end{quote}

\paragraph{Rationale}

提前准备对潜在质疑的反驳，可以增强研究的可信度和说服力

\paragraph{Risk}

反驳可能不够全面或有力

\paragraph{Quotes}
\begin{quote}
Some argue that implicit confidence signals are not reliable
\end{quote}

\section{Method}
\subsection{设计基于隐式 confidence 信号的 DPO 实验框架}
\textit{round 1789354064111}\\
\begin{quote}
Stage: p\_design\_experiment\_plan  ·  Kind: experiment\_plan\_md  ·  Type: experiment\_plan  ·  Round: 1789354064111
Gate decision: sketch  ·  Total score: 5.333333333333333
References: arxiv:2305.16266
\end{quote}

\paragraph{Rationale}

明确实验目标和步骤，为后续实验实施提供清晰的指导

\paragraph{Risk}

实验设计可能过于复杂，难以实施

\paragraph{Quotes}
\begin{quote}
DPO 的 reward hacking 问题可以通过引入额外信号来缓解
\end{quote}

\subsection{添加关于隐式 confidence 信号的最新研究}
\textit{round 1789354064111}\\
\begin{quote}
Stage: p\_design\_experiment\_plan  ·  Kind: experiment\_plan\_md  ·  Type: add\_paper  ·  Round: 1789354064111
Gate decision: sketch  ·  Total score: 6
References: arxiv:2106.03741, arxiv:2206.02790
\end{quote}

\paragraph{Rationale}

确保实验设计基于最新的研究成果，提高实验的创新性和可行性

\paragraph{Risk}

新添加的论文可能与现有研究方向不完全一致

\paragraph{Quotes}
\begin{quote}
隐式 confidence 信号在强化学习中应用广泛
DPO 的改进可以通过引入隐式信号来实现
\end{quote}

\subsection{撰写隐式 confidence 信号在 DPO 中的应用综述}
\textit{round 1789354064111}\\
\begin{quote}
Stage: p\_design\_experiment\_plan  ·  Kind: experiment\_plan\_md  ·  Type: create\_draft  ·  Round: 1789354064111
Gate decision: sketch  ·  Total score: 6
References: arxiv:2305.16266, arxiv:2106.03741, arxiv:2206.02790
\end{quote}

\paragraph{Rationale}

总结现有研究，为实验设计提供理论支持

\paragraph{Risk}

综述可能过于宽泛，缺乏针对性

\paragraph{Quotes}
\begin{quote}
DPO 的 reward hacking 问题
隐式 confidence 信号的定义与应用
DPO 的改进方法
\end{quote}

\subsection{反驳隐式 confidence 信号在 DPO 中应用的潜在问题}
\textit{round 1789354064111}\\
\begin{quote}
Stage: p\_design\_experiment\_plan  ·  Kind: experiment\_plan\_md  ·  Type: rebuttal  ·  Round: 1789354064111
Gate decision: sketch  ·  Total score: 5.666666666666667
References: arxiv:2305.16266
\end{quote}

\paragraph{Rationale}

提前识别并解决潜在问题，增强实验设计的鲁棒性

\paragraph{Risk}

反驳可能过于主观，缺乏实证支持

\paragraph{Quotes}
\begin{quote}
DPO 的 reward hacking 问题可能无法完全通过隐式信号解决
\end{quote}

\subsection{制定隐式 confidence 信号的具体实施步骤}
\textit{round 1789354064111}\\
\begin{quote}
Stage: p\_design\_experiment\_plan  ·  Kind: experiment\_plan\_md  ·  Type: experiment\_plan  ·  Round: 1789354064111
Gate decision: sketch  ·  Total score: 5.333333333333333
References: arxiv:2305.16266
\end{quote}

\paragraph{Rationale}

确保实验的可操作性和可重复性

\paragraph{Risk}

实施步骤可能过于简化，无法涵盖所有细节

\paragraph{Quotes}
\begin{quote}
隐式 confidence 信号的具体实施方法
\end{quote}

\section{Results and Discussion}
\subsection{添加关于隐式 confidence 建模的关键论文}
\textit{round 1789353990210}\\
\begin{quote}
Stage: p\_ideate\_research\_question  ·  Kind: draft\_md  ·  Type: add\_paper  ·  Round: 1789353990210
Gate decision: sketch  ·  Total score: 6
References: 2306.02782, 2212.08068
\end{quote}

\paragraph{Rationale}

需要了解现有的隐式 confidence 建模方法及其在 RLHF 中的应用

\paragraph{Risk}

选择的论文可能与 DPO reward hacking 的具体问题关联度不够

\paragraph{Quotes}
\begin{quote}
隐式 confidence 建模在 RLHF 中可以提供更细粒度的反馈信号
DPO 存在 reward hacking 问题，需要新的隐式信号来缓解
\end{quote}

\subsection{隐式 confidence 信号在 DPO 中的应用综述}
\textit{round 1789353990210}\\
\begin{quote}
Stage: p\_ideate\_research\_question  ·  Kind: draft\_md  ·  Type: create\_draft  ·  Round: 1789353990210
Gate decision: sketch  ·  Total score: 5.333333333333333
References: 2306.02782, 2212.08068, 2210.01241
\end{quote}

\paragraph{Rationale}

整理现有研究，为后续实验设计提供理论基础

\paragraph{Risk}

综述可能过于宽泛，无法直接指导后续研究

\paragraph{Quotes}
\begin{quote}
隐式 confidence 在 RLHF 中的重要性
DPO 的局限性及改进方向
\end{quote}

\subsection{设计基于隐式 confidence 的 DPO 实验方案}
\textit{round 1789353990210}\\
\begin{quote}
Stage: p\_ideate\_research\_question  ·  Kind: draft\_md  ·  Type: experiment\_plan  ·  Round: 1789353990210
Gate decision: sketch  ·  Total score: 5.666666666666667
References: 2212.08068, 2306.02782
\end{quote}

\paragraph{Rationale}

将理论转化为可测试的假设，为验证研究问题做准备

\paragraph{Risk}

实验设计可能过于复杂，难以实现

\paragraph{Quotes}
\begin{quote}
隐式 confidence 可以提供更细粒度的反馈信号
DPO 的 reward hacking 问题需要新的解决方案
\end{quote}

\subsection{反驳 DPO 中隐式 confidence 信号无效的观点}
\textit{round 1789353990210}\\
\begin{quote}
Stage: p\_ideate\_research\_question  ·  Kind: draft\_md  ·  Type: rebuttal  ·  Round: 1789353990210
Gate decision: sketch  ·  Total score: 5.666666666666667
References: 2306.02782
\end{quote}

\paragraph{Rationale}

预判并回应可能的研究质疑

\paragraph{Risk}

反驳可能过于防御性，无法突出研究的创新点

\paragraph{Quotes}
\begin{quote}
隐式 confidence 信号在 RLHF 中已被证明有效
DPO 的 feedback 机制可以通过隐式信号进行改进
\end{quote}

\subsection{DPO 与隐式 confidence 建模的文献综述}
\textit{round 1789353990210}\\
\begin{quote}
Stage: p\_ideate\_research\_question  ·  Kind: draft\_md  ·  Type: literature\_review  ·  Round: 1789353990210
Gate decision: sketch  ·  Total score: 5.333333333333333
References: 2212.08068, 2306.02782, 2210.01241
\end{quote}

\paragraph{Rationale}

系统梳理相关研究领域，为研究问题提供全面背景

\paragraph{Risk}

文献综述可能过于宽泛，无法直接支持研究问题

\paragraph{Quotes}
\begin{quote}
DPO 的主要问题和改进方向
隐式 confidence 在 RLHF 中的应用现状
\end{quote}

\subsection{撰写方法部分的初稿}
\textit{round 1789354102712}\\
\begin{quote}
Stage: p\_write\_paper\_draft  ·  Kind: draft\_md  ·  Type: create\_draft  ·  Round: 1789354102712
Gate decision: sketch  ·  Total score: 4.666666666666667
\end{quote}

\paragraph{Rationale}

方法部分需要详细描述如何利用隐式 confidence 信号来改进 DPO，这将为整个论文奠定基础。

\paragraph{Risk}

方法描述可能不够清晰或完整

\subsection{添加关于 DPO 和 reward hacking 的关键论文}
\textit{round 1789354102712}\\
\begin{quote}
Stage: p\_write\_paper\_draft  ·  Kind: draft\_md  ·  Type: add\_paper  ·  Round: 1789354102712
Gate decision: sketch  ·  Total score: 5.666666666666667
References: 2109.04847, 2203.15860
\end{quote}

\paragraph{Rationale}

确保对 DPO 和 reward hacking 的现有研究有全面的理解，并引用相关工作。

\paragraph{Risk}

可能遗漏一些关键文献

\paragraph{Quotes}
\begin{quote}
DPO 的基本原理
Reward hacking 的常见问题
\end{quote}

\subsection{设计实验来验证隐式 confidence 信号的效果}
\textit{round 1789354102712}\\
\begin{quote}
Stage: p\_write\_paper\_draft  ·  Kind: draft\_md  ·  Type: experiment\_plan  ·  Round: 1789354102712
Gate decision: sketch  ·  Total score: 5.333333333333333
References: 2109.04847
\end{quote}

\paragraph{Rationale}

需要通过实验验证提出的方法是否有效，并量化其对减少 reward hacking 的影响。

\paragraph{Risk}

实验设计可能过于复杂或难以实施

\paragraph{Quotes}
\begin{quote}
实验设计的建议
\end{quote}

\subsection{撰写关于 DPO 和 reward hacking 的文献综述}
\textit{round 1789354102712}\\
\begin{quote}
Stage: p\_write\_paper\_draft  ·  Kind: draft\_md  ·  Type: literature\_review  ·  Round: 1789354102712
Gate decision: sketch  ·  Total score: 5.666666666666667
References: 2109.04847, 2203.15860
\end{quote}

\paragraph{Rationale}

全面回顾现有研究将有助于定位本文的创新点，并展示对领域的深入理解。

\paragraph{Risk}

可能过于宽泛或不够深入

\paragraph{Quotes}
\begin{quote}
DPO 的应用
Reward hacking 的挑战
\end{quote}

\subsection{撰写关于隐式 confidence 信号可能局限性的反驳段落}
\textit{round 1789354102712}\\
\begin{quote}
Stage: p\_write\_paper\_draft  ·  Kind: draft\_md  ·  Type: rebuttal  ·  Round: 1789354102712
Gate decision: sketch  ·  Total score: 6
References: 2109.04847
\end{quote}

\paragraph{Rationale}

预判并回应潜在的批评将增强论文的说服力。

\paragraph{Risk}

可能忽视某些重要的批评

\paragraph{Quotes}
\begin{quote}
DPO 的局限性
\end{quote}

\section{Limitations and Threats to Validity}
\textit{(本节暂无内容 —— rebuttal proposal 产出的自我质疑与答辩。)}

\begin{thebibliography}{99}

\bibitem{arXiv:2106.03741} arXiv:2106.03741, 2021. \url{https://arxiv.org/abs/2106.03741}
\bibitem{arXiv:2109.04847} arXiv:2109.04847, 2021. \url{https://arxiv.org/abs/2109.04847}
\bibitem{arXiv:2109.04882} arXiv:2109.04882, 2021. \url{https://arxiv.org/abs/2109.04882}
\bibitem{arXiv:2203.15860} arXiv:2203.15860, 2022. \url{https://arxiv.org/abs/2203.15860}
\bibitem{arXiv:2206.02790} arXiv:2206.02790, 2022. \url{https://arxiv.org/abs/2206.02790}
\bibitem{arXiv:2210.01241} arXiv:2210.01241, 2022. \url{https://arxiv.org/abs/2210.01241}
\bibitem{arXiv:2210.10757} arXiv:2210.10757, 2022. \url{https://arxiv.org/abs/2210.10757}
\bibitem{arXiv:2212.08068} arXiv:2212.08068, 2022. \url{https://arxiv.org/abs/2212.08068}
\bibitem{arXiv:2212.08073} arXiv:2212.08073, 2022. \url{https://arxiv.org/abs/2212.08073}
\bibitem{arXiv:2302.08582} arXiv:2302.08582, 2023. \url{https://arxiv.org/abs/2302.08582}
\bibitem{arXiv:2305.16266} arXiv:2305.16266, 2023. \url{https://arxiv.org/abs/2305.16266}
\bibitem{arXiv:2305.18290} arXiv:2305.18290, 2023. \url{https://arxiv.org/abs/2305.18290}
\bibitem{arXiv:2306.02782} arXiv:2306.02782, 2023. \url{https://arxiv.org/abs/2306.02782}

\end{thebibliography}

% 注:refs.bib 同样提供 @misc 条目,可用 bibtex 编译替代上面的 thebibliography。
% BibTeX 编译: pdflatex paper && bibtex paper && pdflatex paper && pdflatex paper

\end{document}